% ===========================================================================
\section{Discussion}
\label{sec:discussion}
% ===========================================================================

\subsection{Segment completion changes national accounting}

The common-support experiment identifies a modest increase from the 2006 to the 2019 production
factors. The much larger change appears when gathering, processing, network and post-meter
sources enter the account. For data-scarce countries, applying a newer production factor while
retaining a production-only boundary would therefore capture only a small part of the
Refinement's effect.

The segment-resolved reconstruction reinforces this interpretation. Domestic gas production declines,
but distribution length and customer connections grow. The resulting shift from upstream to
network emissions leaves the national total comparatively stable. A production-scaled time
series would miss that structural change and direct mitigation attention toward a shrinking
share of the system.

\subsection{Two infrastructure questions define the national range}

Distribution practice and post-meter industrial leakage determine most of both the structural
range and the reference-configuration sensitivity. National pipe material, pressure, age and
annual LDAR coverage would establish whether Pakistan meets the lower distribution-factor
criteria. Public utility reports confirm ongoing leak surveys and rehabilitation, but the
quantified material evidence covers only a small part of the network.

The industrial/power factor presents a different problem. Its IPCC lineage and unit conversion
are explicit, and national activity is available by end-use class, but its empirical technology
population comes from an older global default. The full, half-factor and exclusion cases show
the consequence of that transfer. Measurements at power stations, general-industry sites
and captive-power facilities would address the largest end-use components directly.

The reference configuration provides a reproducible application of the public evidence. The
factorial envelopes show how alternative practice, activity-basis and coverage choices affect
that result, while the conditional intervals quantify uncertainty within named configurations.
Together they identify the observations capable of reducing uncertainty instead of compressing
all model choices into one percentage.

\subsection{Utility and official accounts identify reconciliation priorities}

The gas-flow ledger establishes a coherent downstream system from indigenous marketable supply
and RLNG to categorized use and UFG. Its unresolved upstream residual highlights a separate data
need: plant throughput, reinjection, own use, processing shrinkage and operator venting/flaring.
Those quantities would also improve the processing activity basis.

The UFG comparison translates network estimates into a quantity utilities can test. The
reference and lower-practice configurations require markedly different physical shares of the
reported balance, creating a direct role for utility loss audits that separate leakage from
theft, metering and billing effects.

The CRT reconstruction locates the official difference in source representation. Common GWP
conversion and adjacent fiscal years remove metric and period alignment as major explanations;
the public workbook itself is arithmetically coherent. The waterfall instead points to gas
production/gathering activity, processing, distribution, post-meter coverage and the oil--gas
allocation. Reconciliation should therefore begin with the underlying activity bridge and
category boundaries.

\subsection{Satellite observations reveal heterogeneous upstream emissions}

The field analysis supports the broad geography of the upstream inventory while challenging a
simple production-intensity relationship. Observations cover most reported marketable supply,
and MARS source locations are non-random with respect to producing fields. Their modest rank
association with production is consistent with equipment, operating practice and intermittent
events controlling emissions at individual sites.

This finding has two implications. First, production remains useful for allocating an aggregate
upstream estimate when facility activity is incomplete. Second, production cannot substitute
for facility observations when identifying high emitters. Repeated plume coverage, plant and
compressor inventories, and operator-matched field data would turn the present geographic test
into a stronger facility-level evaluation.

\subsection{Implications for data-scarce inventories}

The source-neutral ledger makes future improvements local. Well and field records can replace
national production proxies; plant throughput can resolve processing and disposition; utility
material and pressure classes can select network factors; and facility measurements can replace
transferred post-meter factors. The evidence ranking indicates that national distribution
technology, post-meter measurements and upstream disposition should precede refinements to minor
oil-transport or LNG rows.

More generally, the study shows why factor selection, boundary completion and evaluation should
be treated as separate tasks. A newer factor framework improves an inventory only when its
physical segments are represented, while independent controls are most informative when matched
to the processes they can observe.
